% ===================================================================
%  TABELLA nr. 6 — La chasa dadens, la mobiglia, il nutriment,
%  l'illuminaziun.
%  Traduction 2026 d'apres texto/io, controlee sur texto/fr, texto/ca
%  et texto/oc. Consigne : voir texto/rm/10-tabella-01.tex.
%
%  UNE TRADUCTION DIT AUSSI CE QUE LE LECTEUR N'A PAS CHEZ LUI, ET
%  CETTE PAGE-CI LE DIT PAR LE CHAUFFAGE.
%
%  La planche montre une maison chauffee comme on chauffait dans la
%  plaine : une CHEMINEE dans la chambre et une autre au salon, avec
%  tout son attirail — pelle, pincettes, chenets, garde-cendres,
%  ecran de feu, coffre a buches, pare-etincelles — et, dans la salle
%  a manger seulement, un poele de faience.
%
%  Or une maison de vallee grisonne a renonce a la cheminee il y a
%  des siecles, et n'a garde qu'un seul appareil : LA « PIGNA », le
%  poele de maconnerie, du latin PINEA, la pomme de pin, d'apres la
%  forme de l'epi qui le couronne. C'est le seul mot de chauffage que
%  le romanche possede en propre, et c'est autour de lui que la
%  maison s'organise — la piece qui l'entoure a son nom a elle, la
%  « stiva ».
%
%  LE ROMANCHE DOIT DONC EMPLOYER SON MOT LE PLUS RICHE POUR LA SEULE
%  LIGNE OU LA PLANCHE PARLE DE POELE, ET COMPOSER POUR TOUT LE
%  RESTE : « chavriets » pour les chenets, « stgirm dal fieu » pour
%  l'ecran, « parascintillas » pour le pare-etincelles, « chascha da
%  lain » pour le coffre a buches. Le vocabulaire riche tombe a cote,
%  et les composes couvrent le reste.
%
%  CE N'EST PAS UNE PAUVRETE DE LA LANGUE, C'EST UNE DIFFERENCE DE
%  MAISON. Le tableau 5 avait montre trois couches de metier ; le 6
%  montre que la richesse d'un vocabulaire se mesure a l'objet, pas a
%  la langue : la ou le lecteur romanche a un objet, il a dix mots ;
%  la ou il n'en a pas, il en a un, compose. La colonne
%  luxembourgeoise disait au tableau 10 qu'un manque n'invente rien
%  quand un voisin est la ; ici un manque ne s'invente meme pas un
%  emprunt — il se compose, parce qu'il ne sert qu'a lire.
%
%  RESTITUTION D'UNE PARENTHESE, BLOC t06-c3-08-2. Le fac-simile ido
%  ecrit « « La Hotelo » n° 13) » sans la parenthese ouvrante. Comme
%  aux tableaux 8 et 16, la coquille est du compositeur de l'ido et
%  non du texte romanche : la parenthese est retablie. UNE TRADUCTION
%  N'EST PAS UNE TRANSCRIPTION.
% ===================================================================

%%K t06-apar-1 apar
\VUsaut{6.0mm}
\VUtitre{15.0pt}{60}{TABELLA nr. 6}
\VUsaut{1.4mm}
\VUtitre{11.4pt}{0}{La chasa dadens, la mobiglia, il nutriment,}
\VUsaut{0.4mm}
\VUtitre{11.4pt}{0}{l'illuminaziun.}
\VUsaut{2.2mm}

%%K t06-apar-2 apar
La chasa è finida. L'aug da Gion è installà en sia nova abitaziun, da la quala nus
enconuschain la disposiziun. Uss vulain nus vesair co ch'el ha mobiglià sia chasa.
Nus explorain l'emprim:

%%K t06-tit-1 sub
\VUpk{10.2pt}{40}{La stanza da durmir.}
\VUsaut{3.0mm}

%%K t06-c1-01-1 p
1. --- Nus arrivain a mal mument, perquai che la chombrera è occupada da nettegiar
la \VUgras{chombra}.

%%K t06-c1-02-1 p
2. --- Sin la \VUgras{maisa da toaletta} \textsuperscript{(1)} èn qua e là differents
objects cun ils quals ins sa lava, sa petgna, curtamain ditg ins fa toaletta. Quai
èn il \VUgras{bassin} \textsuperscript{(2)} e la \VUgras{bruncina}
\textsuperscript{(3)}, la \VUgras{biischta da chavels} \textsuperscript{(4)}, il
\VUgras{petgen} \textsuperscript{(5)}, il \VUgras{savun} \textsuperscript{(6)} e la
\VUgras{spugna} \textsuperscript{(7)}, a la fin ina \VUgras{flaschetta}
\textsuperscript{(8)} per l'aua da dents.

%%K t06-c1-03-1 p
3. --- Per terra, avant la maisa da toaletta, qua è la \VUgras{sadella}
\textsuperscript{(9)} per rimnar l'aua sporca.

%%K t06-c1-04-1 p
4. --- Nus vesain en l'\VUgras{anticamera} \textsuperscript{(10)} intginas
\VUgras{hoccas da vestgadira} \textsuperscript{(13)} e duas \VUgras{paraplievgias}
\textsuperscript{(11)} en il \VUgras{portaparaplievgias} \textsuperscript{(12)}.

%%K t06-c1-05-1 p
5. --- Da l'autra vart da la porta è l'\VUgras{armari cun spievel}
\textsuperscript{(14)}, che cuntegna la \VUgras{biancaria} \textsuperscript{(15)}.

%%K t06-c1-06-1 p
6. --- Profitain nus ch'la \VUgras{chombrera} \textsuperscript{(6)} fa il \VUgras{letg}
\textsuperscript{(17)}, per explorar sias differentas parts. Il \VUgras{rom dal letg}
\textsuperscript{(20)} cuntegna in bun \VUgras{matratsch da resorts}
\textsuperscript{(21)}, sin il qual Justina vegn immediatamain a metter in
\VUgras{matratsch da launa} \textsuperscript{(22)}. Lura vegn ella a prender da las
sutgas ils dus \VUgras{leinzuls} \textsuperscript{(23)} e la \VUgras{cuverta}
\textsuperscript{(24)}. Alura vegn ella a metter sin il chau dal letg il
\VUgras{chaschin traversal} \textsuperscript{(25)} ed il \VUgras{chaschin da chau}
\textsuperscript{(26)}, ch'èn cun il \VUgras{plimatsch} \textsuperscript{(27)} sin la
\VUgras{cuverta da letg} \textsuperscript{(32)}. Ella duvra in fasch da plimas per
dispulverar las \VUgras{cortinas} \textsuperscript{(19)} ed il \VUgras{baldachin}
\textsuperscript{(18)}, e qua è ina part da la lavur finida.

%%K t06-c1-07-1 p
7. --- Alura stoppia ella metter in pau en urden la \VUgras{maisetta da letg}
\textsuperscript{(28)}, remontar il \VUgras{svegliarin} \textsuperscript{(29)},
preparar il \VUgras{lampadin da notg} \textsuperscript{(30)}, prender giu la
\VUgras{chandaila} \textsuperscript{(31)} per nettegiar il chandelier.

%%K t06-tit-2 sub
\VUsaut{5.0mm}
\VUpk{10.2pt}{40}{Il chanaster da lavur e las robas dal chamin.}
\VUsaut{3.0mm}

%%K t06-c1-08-1 p
8. --- Il \VUgras{chanaster da lavur} \textsuperscript{(34)} sper il \VUgras{stgabè}
\textsuperscript{(33)} ha era fitg da basegn d'esser mess en urden. Ella stoppia
plijar las \VUgras{bruderias} \textsuperscript{(35)}, serrar en la \VUgras{maisa da
lavur} \textsuperscript{(38)} il \VUgras{digital}, il \VUgras{fil}
\textsuperscript{(36)}, las \VUgras{guaglias} \textsuperscript{(37)} ed il
\VUgras{forsch} \textsuperscript{(39)} e serrar cun la \VUgras{clav}
\textsuperscript{(40)} il cuvertin da la maisa.

%%K t06-c1-09-1 p
9. --- Sin la \VUgras{maisa ad in pe} \textsuperscript{(41)} dal mez vegn Justina a
renovar l'aua da la \VUgras{carafa} \textsuperscript{(42)}. Ella vegn a metter
\VUgras{zucker} en la \VUgras{zuckerera} \textsuperscript{(43)}, a nettegiar la
\VUgras{pinzetta dal zucker} \textsuperscript{(44)} ed a prender giu la
\VUgras{biischta da vestgadira} \textsuperscript{(46)}.

%%K t06-c1-10-1 p
10. --- La vart dretga da la chombra pretenda main lavur, perquai che la
\VUgras{chaise longue} \textsuperscript{(47)} e ses \VUgras{chaschin}
\textsuperscript{(48)} èn a lur dretg plaz, e, perquai ch'i na fa betg fraid,
n'è nagut dislocà tranter las robas dal \VUgras{chamin} \textsuperscript{(49)}. La
\VUgras{pala} \textsuperscript{(50)}, las \VUgras{tanaglias} \textsuperscript{(51)}, ils
\VUgras{chavriets} \textsuperscript{(55)}, la \VUgras{sadella da tschendra}
\textsuperscript{(56)} èn nets; il \VUgras{stgirm dal fieu} \textsuperscript{(54)}
n'è betg vegnì dislocà e la \VUgras{chascha da lain} \textsuperscript{(52)} è bain
pervista da \VUgras{lennas} \textsuperscript{(53)}.

%%K t06-noto-1 noto
\VUnotes{143.2mm}{%
(*) Babo = uffantin, l'E. baby, F. bébé, I. bambino.
\par
%%K t06-noto-2 noto
(*) Lambrequino = F. lambrequin, S. lambrequines, Port. lambrequins, era D. E.
lambrequins, pia D. E. F. P. S.%
}

%%K t06-c1-11-1 p
11. --- Tuttina stoppia ella dispulverar la \VUgras{pendula} \textsuperscript{(57)},
ils \VUgras{candelabers} \textsuperscript{(58)} ed ils \VUgras{vaschels da giaglioffa}
\textsuperscript{(59)}, a la fin sfruschar il \VUgras{spievel} \textsuperscript{(60)}.

%%K t06-c1-12-1 p
12. --- Quant a la \VUgras{cuna} \textsuperscript{(61)} da l'uffantin (dal babo (*)),
ella è gia pronta.

%%K t06-tit-3 sub
\VUsaut{5.0mm}
\VUpk{10.2pt}{40}{Il salun.}
\VUsaut{3.0mm}

%%K t06-c2-01-1 p
1. --- Uss qua è il salun. Quai è ina fitg bella \VUgras{chombra}. Il chamin, ch'è
en l'chantun dretg, è ornà cun ina delicata \VUgras{statuetta} \textsuperscript{(1)},
cun dus bels \VUgras{vaschs} \textsuperscript{(3)} e drappà cun in \VUgras{lambrequin}
\textsuperscript{(2)} da velur (*).

%%K t06-c2-02-1 p
2. --- Per impedir ch'las scintillas dal fieughin brischian il tapet, ha ins mess
avant il fieughin in \VUgras{parascintillas} \textsuperscript{(4)} da rom.

%%K t06-c2-03-1 p
3. --- Dasperas è avert in elegant \VUgras{secretari} \textsuperscript{(5)} da dunna.
Sin el scriva la \VUgras{patruna da chasa} \textsuperscript{(23)} sias brevs.

%%K t06-c2-04-1 p
4. --- La fanestra è ornada cun \VUgras{cortinas da fanestra} \textsuperscript{(6)}
cun \VUgras{ligiams} \textsuperscript{(8)} tacads als \VUgras{pommels}
\textsuperscript{(7)} e cun grondas cortinas da stoffa \VUgras{drappadas dad ensi}
\textsuperscript{(9)}.

%%K t06-c2-05-1 p
5. --- En la giaviada da la fanestra cuntegna in \VUgras{portavasch}
\textsuperscript{(11)} ina verda \VUgras{planta} \textsuperscript{(10)}, ina
splendida palma, da la quala ils fegls sa stendan quasi fin a la \VUgras{lampa da
petroli} \textsuperscript{(12)}.

%%K t06-c2-06-1 p
6. --- Il \VUgras{paraglisch} \textsuperscript{(13)} da la lampa è vegnì fatg da
Maria, la charmanta \VUgras{giuvna} \textsuperscript{(14)} ch'è al piano
\textsuperscript{(15)}. Sesend fitg dretga sin il \VUgras{stgabè}
\textsuperscript{(16)}, studegia ella in toc da musica. Ella è fitg abila e
premurusa, sco che mussa ses \VUgras{armari da musica} \textsuperscript{(17)}, ch'è
perfectamain en urden.

%%K t06-tit-4 sub
\VUsaut{5.0mm}
\VUpk{10.2pt}{40}{Il mat nauscha.}
\VUsaut{3.0mm}

%%K t06-c2-07-1 p
7. --- Ses frar, Paul, tschertga in \VUgras{cudesch} \textsuperscript{(19)} en la
\VUgras{biblioteca} \textsuperscript{(18)}. El è in \VUgras{giuven}
\textsuperscript{(20)}, malquiet e nunsuffribel. Sa tacond a la biblioteca per
cuntanscher il cudesch ch'è memia aut, ristga el da far crudar il \VUgras{bust}
\textsuperscript{(21)} ch'è sisum.

%%K t06-c2-08-1 p
8. --- El profitescha, per far quella stupidezza, ch'sia mamma è occupada a
discurrer cun in'amia, ina \VUgras{dama} \textsuperscript{(22)} ch'è vegnida a passar
intgins dis en sia chasa. La mamma sesa sin in \VUgras{canapé}
\textsuperscript{(24)} sper il \VUgras{sutgun} \textsuperscript{(25)}, e sia amia è sin
la \VUgras{sutga} \textsuperscript{(27)}, da la quala ins vesa mo la
\VUgras{spatlera} \textsuperscript{(26)}.

%%K t06-c2-09-1 p
9. --- Il grond \VUgras{rom} \textsuperscript{(30)}, che penda vi da la paraid,
cuntegna il \VUgras{portret} \textsuperscript{(29)} d'ina parenta. En l'\VUgras{album}
\textsuperscript{(32)} mess sin la maisa èn radunadas las \VUgras{fotografias}
\textsuperscript{(33)} dals auters members da la famiglia. Ils uffants han gist
sfegliettà en el, perquai ch'las \VUgras{sutgas} \textsuperscript{(31)} èn anc
radunadas enturn la maisa.

%%K t06-c2-10-1 p
10. --- Il \VUgras{vasch chinais} \textsuperscript{(34)} da porzelain, ch'è sper il
\VUgras{curtè da palpiri} \textsuperscript{(35)}, è vegnì regalà a Maria per sia
festa, ed il \VUgras{paravent} \textsuperscript{(36)} ch'orna il chantun da la
chombra è vegnì purtà da la China da ses aug.

%%K t06-c2-11-1 p
11. --- Allegrà da las bellissimas \VUgras{picturas} \textsuperscript{(37)} tacadas
vi da la \VUgras{paraid} \textsuperscript{(42)}, illuminà dal \VUgras{luster}
\textsuperscript{(38)}, dal qual las \VUgras{lampas electricas}
\textsuperscript{(39)} splenduran allegramain la saira, para quel salun fitg
agreabel. Il mol \VUgras{tapet} \textsuperscript{(40)} stendì sin il parquet e la
\VUgras{clacca electrica} \textsuperscript{(41)}, uschè cumadaivla, cumpletteschan ses
confort.

%%K t06-tit-5 sub
\VUsaut{3.6mm}
\VUpk{10.2pt}{40}{La stanza da magliar.}
\VUsaut{3.0mm}

%%K t06-c3-01-1 p
1. --- La campana da la tschaina ha sunà. Tut la famiglia, cun excepziun da Maria,
è radunada enturn la maisa. La stanza da magliar è agreablamain stgaudada da la
excellenta \VUgras{pigna} \textsuperscript{(1)} da fajenza, cun in sutgl \VUgras{tub}
\textsuperscript{(2)}.

%%K t06-c3-02-1 p
2. --- Quella chombra n'ha betg bler mobiglia. Lung il mir penda, sur la
\VUgras{banchetta} \textsuperscript{(4)}, ina \VUgras{funtanella} \textsuperscript{(3)}
ornamentala. Fitg dasperas è la \VUgras{credenza} \textsuperscript{(5)}, sin la quala
ins metta las spaisas fraidas, ils plats da dessert e las differentas robas da
maisa ch'ins n'ha betg immediatamain da basegn.

%%K t06-c3-03-1 p
3. --- Uschia vesain nus sin l'assa da sura in \VUgras{pulschin rustì}
\textsuperscript{(7)}, in \VUgras{pesch} \textsuperscript{(8)} ed ina \VUgras{cuppa}
\textsuperscript{(10)} cun \VUgras{fritgas} \textsuperscript{(9)}. Sutvart qua è il
\VUgras{chaschiel} \textsuperscript{(11)}, il \VUgras{paun} \textsuperscript{(12)}, il
\VUgras{portaflaschettas} \textsuperscript{(13)}, che cuntegna tut quai ch'ins ha da
basegn per cundir la salata: l'ieli, l'aschieu, il \VUgras{sal}
\textsuperscript{(14)} ed il \VUgras{paiver} \textsuperscript{(15)}. A la fin è sin
l'assa da sut ina \VUgras{tuorta} \textsuperscript{(17)} sper la \VUgras{senavrera}
\textsuperscript{(16)}.

%%K t06-c3-04-1 p
4. --- L'ura da la stanza da magliar, ch'ins numna \VUgras{ura da paraid}
\textsuperscript{(20)}, ha ina furma fitg speziala. Ella è messa en in rom da lain,
che cuntegna plinavant in \VUgras{barometer} \textsuperscript{(19)} ed in
\VUgras{termometer} \textsuperscript{(18)}.

%%K t06-tit-6 sub
\VUsaut{4.6mm}
\VUpk{10.2pt}{40}{Il servir da la tschaina.}
\VUsaut{3.0mm}

%%K t06-c3-05-1 p
5. --- Vi da tut ils mirs da la chombra èn applitgads \VUgras{panels da lain}
\textsuperscript{(21)} da ruver incerà. Ina \VUgras{portiera} \textsuperscript{(22)}
da stoffa serra avunda bain la porta che maina en la veranda. Là vegn la famiglia
ad ir a beiver il caffè suenter la tschaina. Gia èn las \VUgras{scadiolas}
\textsuperscript{(24)} e las \VUgras{sutscadiolas} \textsuperscript{(25)} prontas sin la
\VUgras{mobiglia dals plats} \textsuperscript{(23)}.

%%K t06-c3-06-1 p
6. --- Sur la maisa è fixada al plafun ina \VUgras{suspensiun} \textsuperscript{(26)},
da la quala la lampa fitg splendurusa illuminescha la saira tut la stanza da
magliar.

%%K t06-c3-07-1 p
7. --- Uss explorain nus la \VUgras{maisa da magliar} \textsuperscript{(27)} sezza.
Cur ch'la famiglia è entrada, era la maisa parvida. Sin la \VUgras{tavatscha}
\textsuperscript{(28)} era mess, al plaz da mintga cumpogn da maisa, in
\VUgras{plat} \textsuperscript{(33)}, ina \VUgras{servietta} \textsuperscript{(29)},
ina \VUgras{chadagna} \textsuperscript{(34)}, ina \VUgras{furtgetta}
\textsuperscript{(35)}, in \VUgras{curtè} \textsuperscript{(36)} ed in
\VUgras{magiel} \textsuperscript{(32)}. I era era \VUgras{buttiglias}
\textsuperscript{(30)} per il vin ed ina \VUgras{carafa} \textsuperscript{(31)} per
l'aua.

%%K t06-c3-08-1 p
8. --- Il \VUgras{servitur} \textsuperscript{(39)} ha l'emprim purtà la
\VUgras{suppiera} \textsuperscript{(37)}, en la quala fima anc in pau schuppa. Uss
serva el l'emprim \VUgras{plat} \textsuperscript{(38)}, in plat da charn. Lura vegn el
a preschentar quels che nus avain gia numnà, e, a la fin, suenter il \VUgras{plat da
dessert}, vegn el a profitar ch'ses patruns èn ids en la veranda, per prender davent
las robas da maisa e per metter danovamain en urden la stanza da magliar.

%%K t06-c3-08-2 p
\textit{Nota.} --- \textit{Ils pleds da diever cuminaivel che pertutgan il
nutriment vegnan qua inditgads fitg curtamain. La glista vegn cumplettada sin las
duas tabellas «L'hotel» (nr.~13) ed «Il martgà» (nr.~15).}

%%K t06-tit-7 sub
\VUsaut{4.6mm}
\VUpk{10.2pt}{40}{La cuschina.}
\VUsaut{3.0mm}

%%K t06-c4-01-1 p
1. --- En la cuschina, en la quala nus guardain tras la porta mez averta, vesain nus
in pittoresc disurden. Avant il \VUgras{fieughin} \textsuperscript{(1)}, en il qual
crepitescha in \VUgras{fieu} \textsuperscript{(3)}, è installà in \VUgras{virarost}
\textsuperscript{(5)}. In bel \VUgras{rustì} \textsuperscript{(7)} è \VUgras{mess sin
il spid} \textsuperscript{(6)} e finescha da cuschinar.

%%K t06-c4-02-1 p
2. --- Il \VUgras{sufflet} \textsuperscript{(8)} e la \VUgras{pala da carvun}
\textsuperscript{(9)}, ch'ins ha gist duvrà, èn restads per terra sco era las
\VUgras{lennas} \textsuperscript{(10)}.

%%K t06-c4-03-1 p
3. --- Il sisum dal chamin è en urden; la \VUgras{lantiarna} \textsuperscript{(11)}, la
\VUgras{chaschetta da zulprins} \textsuperscript{(13)}, en la quala èn ils
\VUgras{zulprins} \textsuperscript{(12)}, il \VUgras{chandelier}
\textsuperscript{(14)} ed il \VUgras{portachandaila} \textsuperscript{(15)} cun ina
\VUgras{chandaila} \textsuperscript{(16)} èn a lur plaz. Ma la gronda \VUgras{sadella
da carvun} \textsuperscript{(18)}, plaina da \VUgras{carvun da terra}
\textsuperscript{(17)}, ch'la \VUgras{cuschinunza} \textsuperscript{(23)} ha gist
emplenì en il chaler, è restada en il mez da la cuschina.

%%K t06-c4-04-1 p
4. --- Propi, la paupra dunna sto sa prescher per ch'il gentar saja pront a temp.
Uss è ella sper il \VUgras{fuorn da cuschina} \textsuperscript{(19)} e stira ina
\VUgras{sosa} \textsuperscript{(24)} ch'ins na dastga betg laschar brischar.

%%K t06-c4-05-1 p
5. --- En il \VUgras{fuorn} \textsuperscript{(20)} cuscha ina tuorta ch'ella ha
preparà per la marenda dals uffants. Sur il fuorn da cuschina pendan il
\VUgras{chadagnun} \textsuperscript{(21)} ed il \VUgras{sbrumadur}
\textsuperscript{(22)}.

%%K t06-tit-8 sub
\VUsaut{4.0mm}
\VUpk{10.2pt}{40}{Ils vaschels.}
\VUsaut{3.0mm}

%%K t06-c4-06-1 p
6. --- La \VUgras{battaria da cuschina} \textsuperscript{(25)}, tacada en il fund da
la chombra, è fitg netta. Las bellas \VUgras{casserolas} \textsuperscript{(26)} da
rom, il \VUgras{vasch da latg} \textsuperscript{(27)}, la \VUgras{criva}
\textsuperscript{(28)} e la \VUgras{raspa} \textsuperscript{(29)} èn vegnids nettegiads
premurusamain.

%%K t06-c4-07-1 p
7. --- La \VUgras{chascha da sal} \textsuperscript{(30)} è messa sin l'armari dals
plats sper la \VUgras{tekanna} \textsuperscript{(31)} e la \VUgras{buttiglia d'ieli}
\textsuperscript{(32)}. La \VUgras{tirunca} \textsuperscript{(33)} cuntegna
probablamain las posadas ed ils curtels da cuschina.

%%K t06-c4-08-1 p
8. --- La \VUgras{scua} \textsuperscript{(34)} ed il \VUgras{fasch da plimas}
\textsuperscript{(35)} èn en in chantun. La cuschinunza ha gist duvrà il \VUgras{bloc
da lain} \textsuperscript{(36)}, ella l'ha laschà sper la fanestra.

%%K t06-c4-09-1 p
9. --- In \VUgras{sfruschamaun} \textsuperscript{(37)} penda vi dal mir, sper
l'\VUgras{imbut} \textsuperscript{(38)}. En quella part da la cuschina èn messas ensi
la \VUgras{grilla} \textsuperscript{(39)}, il \VUgras{curtè da tagliar}
\textsuperscript{(40)} e l'\VUgras{assa da tagliar} \textsuperscript{(41)}. Lura, sur
il curtè da tagliar, sin ina \VUgras{assa} \textsuperscript{(42)}, èn
\VUgras{vaschs} \textsuperscript{(43)}, la \VUgras{caldera} \textsuperscript{(44)} cun
ses \VUgras{cuvertin} \textsuperscript{(45)} e la \VUgras{marmita}
\textsuperscript{(46)}. La \VUgras{padella} \textsuperscript{(47)} penda dasperas.

%%K t06-c4-10-1 p
10. --- Mo il \VUgras{rubinet} \textsuperscript{(48)} da rom dat in pau da glisch en
quest chantun. Il \VUgras{lavandin} \textsuperscript{(49)}, sin il qual è mess il
gross \VUgras{brunc} \textsuperscript{(50)}, è probablamain fitg cumadaivel cun ses
pitschen armari, en il qual ins conserva il \VUgras{barigl d'aschieu}
\textsuperscript{(51)} pervis cun ses \VUgras{rubinet} \textsuperscript{(52)}, la
\VUgras{chascha dals rests} \textsuperscript{(53)} ed ils \VUgras{plats nunnets}
\textsuperscript{(54)}.

%%K t06-tit-9 sub
\VUsaut{4.0mm}
\VUpk{10.2pt}{40}{Auters objects mess en la cuschina.}
\VUsaut{3.0mm}

%%K t06-c4-11-1 p
11. --- La \VUgras{tina} \textsuperscript{(55)}, en la quala ins lava las
\VUgras{verduras} \textsuperscript{(58)}, e la \VUgras{caldera} \textsuperscript{(56)}
da fier fusì messa sin il \VUgras{palantschieu da tegels} \textsuperscript{(57)} han
probablamain era lur plaz en quell'armari.

%%K t06-c4-12-1 p
12. --- La cuschinunza n'ha segir betg mancanza da fuorns, perquai che qua è
plinavant, sin il chantun da la maisa, in \VUgras{fuornet da gas}
\textsuperscript{(59)}, sin il qual sa stgauda aua en ina pitschna casserola. Il
\VUgras{vapur} \textsuperscript{(60)} che va ora d'ella ans mussa ch'ella bula
probablamain. Presumablamain vegn quell'aua duvrada per il caffè, perquai che qua èn
sin l'autra fin da la maisa la \VUgras{cafetiera} \textsuperscript{(64)} ed il
\VUgras{malin da caffè} \textsuperscript{(63)}.

%%K t06-c4-13-1 p
13. --- Il fuornet è collià cun il \VUgras{brischader da gas} \textsuperscript{(62)}
tras in lung \VUgras{tubet da gumma} \textsuperscript{(61)}.

%%K t06-c4-14-1 p
14. --- La \VUgras{servienta a di} \textsuperscript{(66)}, che vegn a gidar la
cuschinunza, prepara las verduras per la tschaina. Cur ch'ellas vegnan a esser
mundadas e lavadas, vegn ella a las metter en il \VUgras{bassin}
\textsuperscript{(65)}, spetgond l'ura da las cuschinar.

%%K t06-tit-10 sub
\VUsaut{4.0mm}
{\centering\fontsize{10.2pt}{10.2pt}\selectfont\textls[40]{\scshape La stanza da bogn.} \textit{(Chombra da bogn.)}\par}
\VUsaut{3.0mm}

%%K t06-c5-01-1 p
1. --- Ans resta anc mo ina indiscreziun da far per enconuscher las chombras
principalas da l'abitaziun: vesair l'installaziun da la stanza da bogn. Quai na
pretenda betg ditg, perquai ch'ella n'è betg gronda, e nus savain svelt ch'la
\VUgras{tina da bogn} \textsuperscript{(1)} ha in bun \VUgras{stgaudabogn}
\textsuperscript{(2)}, ch'la \VUgras{conca da savun} \textsuperscript{(3)} è
cuntanschibla per il maun da quel che sa bogna e ch'il \VUgras{portapezzas}
\textsuperscript{(5)} fa segir sitgentar facilmain las \VUgras{pezzas da toaletta}
\textsuperscript{(4)}.

%%K t06-c5-02-1 p
2. --- Quella chombra ha sias paraids cuvridas cun \VUgras{plattas da fajenza}
\textsuperscript{(6)}, che han permess d'installar ina \VUgras{duscha}
\textsuperscript{(7)} senza avair tema che l'aua, che sprizza enavos durant il
diever, donnegia ils mirs. In \VUgras{bassinet} \textsuperscript{(8)} da zinc, che
serva per ils bogns da peis, cumplettescha la mobiglia.

\VUsaut{6.0mm}
\VUfilet{22mm}
